% !PRINTED 137
\weberpage{137}
%<<CONT
die wenigen darunter enthaltenen nicht metacyklischen und besonders die einfachen nicht cyklischen Gruppen.

Es ist bis jetzt nicht gelungen, die Gradzahlen der einfachen Gruppen in einem bestimmten Gesetze zusammenzufassen, und man hat sich begnügt, einerseits mit Hülfe der oben entwickelten Sätze, andererseits durch besondere Methoden alle einfachen Gruppen bis zu gewissen Grenzen der Gradzahlen zu ermitteln. \gesperrt{Hölder} hat diese Untersuchungen für die Gradzahlen bis 200 und \gesperrt{Cole} für die Gradzahlen bis 500 durchgeführt\footnote{\gesperrt{Hölder}, Mathematische Annalen, Bd.\,40. \gesperrt{Cole}, American Journal, Vol.\,14.}. Von einfachen nicht cyklischen Gruppen haben sich dabei nur die auch schon aus anderen Untersuchungen bekannten von den Graden 60, 168, 360 ergeben. Es ist ferner noch eine einfache Gruppe vom Grade 660 bekannt, und neuerdings haben \gesperrt{Cole} und \gesperrt{Moore} noch auf eine einfache Gruppe vom Grade 504 aufmerksam gemacht\footnote{Bulletin of the New York mathem. society, Oct.\,1893.}. Diese Untersuchungen werden mit dem Wachsen der Gradzahlen sehr mühsam. Wir wollen uns hier auf die Betrachtung der Gradzahlen des ersten Hunderts beschränken.

Hier erweist sich nach den drei allgemeinen Sätzen von \gesperrt{Sylow} und \gesperrt{Frobenius} die Mehrzahl der Gruppen als metacyklisch, und es bleiben nur die Gradzahlen

\[ 36,\ 60,\ 72,\ 84,\ 90,\ 100 \]

übrig. Dass aber eine Gruppe $36^{\mathrm{sten}}$ Grades nicht einfach sein kann, ergiebt sich nach dem \gesperrt{Sylow}'schen Satze. Denn eine solche Gruppe müsste einen Theiler $9^{\mathrm{ten}}$ Grades haben und müsste also durch die Permutationen von vier Ziffern darstellbar sein. Das ist aber unmöglich, weil es nur 24 Permutationen von vier Ziffern giebt. Eine Gruppe vom Grade $72 = 8 . 9$ muss nach \S.\,30, II., da 8 nicht $\equiv 1 \pmod{3}$ ist, einen Theiler vom Grade 18 enthalten, und wäre also, wenn sie einfach wäre, ebenfalls durch die Permutationen von vier Ziffern darstellbar, was noch weniger möglich ist. Dass einfache Gruppen von den Graden 84 und 100 nicht existiren können, ergiebt sich gleichfalls sofort aus den \gesperrt{Sylow}'schen Sätzen, da, wenn wir die Gruppen $7^{\mathrm{ten}}$ oder $25^{\mathrm{sten}}$ Grades aussondern, kein Theiler übrig bleibt, der nach dem Modul 7 oder 5 mit 1 congruent ist. Dass alle diese Gradzahlen nur metacyklischen Gruppen angehören können, folgt dann nach \S.\,6, 2. ohne Weiteres daraus, dass die Gruppen, deren Grade
%CONT>>
